\documentclass[psamsfonts]{amsart}

%-------Packages---------
\usepackage{amssymb,amsfonts}
\usepackage[all,arc]{xy}
\usepackage{enumerate}
\usepackage{mathrsfs}
\usepackage{graphicx}
\usepackage{hyperref}
\hypersetup{
colorlinks=True,
allcolors=black
}

%--------Theorem Environments--------
%theoremstyle{plain} --- default
\newtheorem{thm}{Theorem}[section]
\newtheorem{cor}[thm]{Corollary}
\newtheorem{prop}[thm]{Proposition}
\newtheorem{lem}[thm]{Lemma}
\newtheorem{conj}[thm]{Conjecture}
\newtheorem{quest}[thm]{Question}

\theoremstyle{definition}
\newtheorem{defn}[thm]{Definition}
\newtheorem{defns}[thm]{Definitions}
\newtheorem{con}[thm]{Construction}
\newtheorem{exmp}[thm]{Example}
\newtheorem{exmps}[thm]{Examples}
\newtheorem{notn}[thm]{Notation}
\newtheorem{notns}[thm]{Notations}
\newtheorem{addm}[thm]{Addendum}
\newtheorem{exer}[thm]{Exercise}

\theoremstyle{remark}
\newtheorem{rem}[thm]{Remark}
\newtheorem{rems}[thm]{Remarks}
\newtheorem{warn}[thm]{Warning}
\newtheorem{sch}[thm]{Scholium}

\makeatletter
\let\c@equation\c@thm
\makeatother
\numberwithin{equation}{section}

\bibliographystyle{plain}

%--------Meta Data: Fill in your info------
%Title Portfolio BC I/II Portfolio \\ Semester Year
\title{Template for Portfolio}

\author{First Author and Second Author}

\date{\today}

\begin{document}

% Delete the abstract
\begin{abstract}

This is a sample latex document with emphasis on using math mode and
equation environments.  You should use it as a template for your portfolio.
This document was adapted from J. P. May's template / advice for writing REU papers, taken from \href{https://math.uchicago.edu/~may/REU2020/}{https://math.uchicago.edu/~may/REU2020/}.



\end{abstract}

\maketitle

\tableofcontents
%Your portfolio will have a section for each unit. You should include 2 proofs from each unit. 
\section{Writing}  Before getting to latex comments, I'll say a few words about writing, venting from many years of hard experience. 

It is important that what you write is something you would actually like to read.  There should be no overuse of symbols.
Clusters of symbols can be barbaric.  There are subjects whose literature is festooned with sentences written almost entirely with symbols like $\exists$, $\forall$.  Don't do that. Never use a symbol for a verb.  Never start a sentence with a symbol.  

Grammar should be correct.  Never start a sentence with ``Where ...".   
Mismatches of singular and plural are excruciatingly painful, utterly abhorrent.  You cannot write
``Let ..., then ..."  That is what is called a run-on sentence. You must write ``Let ... .  Then ... ."
Alternatively, ``If ..., then ..."  works just fine.

Avoid words or phrases like Simply, Obviously, Just, ..., It is easy to see.  They serve only to intimidate or to browbeat the reader into acquiescence. 


References to numbered statements are best in the form Theorem 2.3, not just 2.3.
Equations should be referred to as (2.3), NOT 2.3 or equation 2.3 or equation (2.3).


Bad writing makes for unpleasant papers, no matter how good the material.


\section{What does the table of contents command do?}

The table of contents command will automatically make a contents.

\section{How do the environment commands work?}

``Environments'' are commands that are given using the \verb|\begin{}|
  and \verb|\end{}| syntax. In the preamble, you can see we've defined
  a bunch of theorem-type environments.  For example, to get a definition, 
you type:

\begin{defn}  This is how to define a definition.
\end{defn}

And for a theorem and its proof you type:

\begin{thm}
This is the statement of a theorem.
\end{thm}
\begin{proof}
And this shows that the statement is correct.
\end{proof}

Note that the numbering is taken care of automatically, and that we've predefined a bunch of these sorts of environments in the header.  These take give environments for  lemmas, corollaries, etc. 

Another useful kind of enviroment is the equation environment.  Equations
get numbered in sequence with statements, as for example

\begin{equation}  e = mc^2
\end{equation}

Note that if you do not want a numbered equation, you can use the
environment ``equation*''
 like so:

\begin{equation*}
e=mc^2
\end{equation*}
But a quicker way to tex the same command and get the same displayed result is:
\[  e=mc^2 \]

There are plenty of other equation-type enviroments that allow you to
align several equations and such like things. The AMS's guide \cite{amsshort} is a
good place to start with these.

You can also typeset math directly in a paragraph by placing it within
dollar signs.  This is called ``math mode.''  For example: Let $e$
be energy, $m$ be momentum and $c$ be the speed of light.  Then
Einstein's famous equation says that $e=mc^2$.  This is useful, but
remember that it is harder to read inline math than displayed math. 

Remember that letters get put in a different font in math mode, so
whenever you are referencing a mathematical object you should always
put it in dollar signs.  For example, $f$ is a function, but f is just
a random letter.

Both \cite{notsoshort} and \cite{amsshort} have good lists of other symbols you can use in math mode.
These include greek letters ($\alpha, \beta, \Gamma, \Delta$),
operators ($\otimes, +, \sum$) and much more ($\leq, \diamond$).

\newpage 
\section{Figures and tables}

Here is an example of a table:

\begin{table}[h]
    \centering
    \begin{tabular}{|c|c|} \hline
        $x$ & $y$  \\  \hline
         1 & 1 \\  \hline
         2 & 4 \\  \hline
         3 & 9  \\ \hline 
    \end{tabular}
    \caption{The first three perfect squares.}
    \label{tab:perfect_squares}
\end{table}

The easiest way to include a graphic is to take a screenshot (or draw it by hand and take a picture with your phone), save it as a png or jpg, then upload it to overleaf. To keep things organized, it is a good idea to make a folder called "images" to keep all images. You can display your image using the commands below. You may have to play with the setting for "scale" to adjust the size.

\begin{figure}[h]
    \centering
    \includegraphics[scale=0.3]{images/sample_diagram.png}
    \caption{A sample diagram.}
    \label{fig:sample_diagram}
\end{figure}

\section*{Acknowledgments}  You should thank anyone who you consulted with. For example, ``I thank [another student] for helping me understand [something or other]"

\section{bibliography}  Here is a sample bibliography. You probably won't need it for your portfolio. You do not need to cite the notes from class. More generally, a bibliography should list all sources that you have used and referenced. And you should reference anything you use.   Especially if you quote any result without proof, you MUST
give a reference.   And never ever should you copy material directly or more or less directly, from a source.

\begin{thebibliography}{9}

\bibitem{ams} http://www.ams.org/publications/authors/tex/amslatex

\bibitem{amsshort}
Michael Downes.
Short Math Guide for \LaTeX.
http://tex.loria.fr/general/downes-short-math-guide.pdf

\bibitem{May}
J. P. May.
A Concise Course in Algebraic Topology.
University of Chicago Press. 1999. 

\bibitem{notsoshort}
Tobias Oekiter, Hubert Partl, Irene Hyna and Elisabeth Schlegl.
The Not So Short Introduction to \LaTeX 2e.
https://tobi.oetiker.ch/lshort/lshort.pdf

\end{thebibliography}

\end{document}

